\itemhead{习题2.1 · 第12题（教材p.44）}{XJ-B1-C02-HW-P044-Y1420}
\label{XJ-B1-C02-HW-P044-Y1420}
{\footnotesize 湘教版教材p.44 / PDF第51页 / 习题2.1 / 第12题}
\field{原题}{某公司计划在本地租地建仓库，已知每月土地费用与仓库到车站的距离成反比，而每月货物的运输费用与仓库到车站的距离成正比。经测算，若在距离车站 $10\,\mathrm{km}$ 处建仓库，则每月的土地费用和运输费用分别为 2 万元和 8 万元。那么要使两项费用之和最小，仓库应建在离车站多远处？}
\field{参考答案}{设距离为 $x>0$ 千米，费用 $20/x+0.8x\ge8$ 万元；等号 $x=5$。应建在距车站5千米处。}
{\small \field{关联知识}{I03 基本不等式、等号与最值建模（教材帮p.52、53、54 / PDF61、62、63）；F03 函数求值、解析式与函数建模（教材帮p.80、87、88、89、90、131 / PDF89、96、97、98、99、140）}
\field{方法与易错点}{列正数条件与固定量，构造基本不等式，解等号条件并代回实际限制。 易错：把一个下界误当最小值；多个等号条件不能同时实现；漏单位和整数限制。}
\field{答案依据}{答案册相关题组：PDF第12、13页（册内页11、12）；按当前题库原题逐项重算，不按旧版题号机械对套。}}
